% ============================================================================
% CONCLUSION - REVISED
% ============================================================================

\section{Conclusion}

We have introduced PRISM, a memory-augmented neuro-symbolic framework that addresses two fundamental limitations of existing neural-symbolic approaches to CSP solving: ineffective repair mechanisms and absence of cross-problem experience accumulation.

\subsection{Key Contributions}

PRISM's design rests on a simple but powerful insight: because every inference step in CSP solving can be mechanically verified via SMT solvers, experience accumulation can be formal rather than heuristic. This leads to four concrete contributions:

\begin{enumerate}
    \item \textbf{Formally Verified Paradigm Library:} We introduce the first paradigm formalism where solving strategies are encoded as Z3-verifiable five-tuples (trigger conditions, operations, pre/post-conditions, scope). This enables mechanical certification of experience quality, distinguishing PRISM from all prior memory-based methods relying on unverified natural language.

    \item \textbf{Automated Paradigm Distillation:} We design a fully unsupervised offline pipeline that extracts ~35 high-confidence solving paradigms from 1500+ training trajectories via KDP identification, cross-trajectory clustering, LLM abstraction, and triple Z3 verification—without human annotation.

    \item \textbf{Adaptive Repair Through Typed Memory:} We introduce a repair trajectory memory system that records UNSAT cores (formally localized contradictions), detects stagnation and loops via Jaccard and cosine similarity, and triggers a four-level strategy escalation protocol. This mechanically breaks repair cycles—a persistent failure mode in prior neuro-symbolic systems.

    \item \textbf{Empirical Validation and Analysis:} Comprehensive evaluation on ZebraLogic demonstrates 14.4pp improvements over prior neuro-symbolic methods, with gains increasing monotonically with problem difficulty (14.1pp on 3×5, 13.5pp on 6×6), validating our thesis that experience accumulation provides exponentially increasing value for harder instances.
\end{enumerate}

\subsection{Key Findings}

Across 800 test puzzles from ZebraLogic:

\begin{itemize}
    \item \textbf{Performance:} PRISM achieves 50.4\% average accuracy, a 40\% relative improvement over the prior neuro-symbolic approach (36.0\%).

    \item \textbf{Repair Efficiency:} Repair rounds decrease from 3.9 (prior work) to 2.1 (PRISM)—a 45\% reduction. Stagnation detection fires on 22\% of repair cases, with 72\% ultimately succeeding after strategy escalation.

    \item \textbf{Paradigm Quality:} The 40:1 compression ratio (1500 trajectories → 35 paradigms) with 94\% hit rate shows that abstraction captures genuine solving patterns. Paradigms maintain reasonable transfer across scales (4--6pp degradation) and partially transfer across domains (6pp accuracy gain on Knights-and-Knaves).

    \item \textbf{Component Synergy:} Paradigm library alone yields 44.4\% (8.4pp gain); repair memory alone yields 41.2\% (5.2pp gain); combined they achieve 50.4\%—demonstrating positive interaction between components.
\end{itemize}

\subsection{Limitations and Future Work}

\textbf{Scope:} This work focuses on logical grid puzzles (Zebra/Knights-and-Knaves), which are representative CSP instances but lack some properties of general CSPs (e.g., continuous domains, general constraint functions). Future work should explore paradigm distillation for broader CSP classes (scheduling, resource allocation, graph coloring, etc.).

\textbf{Scalability:} While PRISM shows gains up to 6×6 scales, extending to significantly larger problems (10×10 grids, 100+ variables) requires investigation into offline efficiency (clustering may not scale) and online efficiency (paradigm retrieval and consistency checking costs).

\textbf{Paradigm Discovery:} The current clustering-based approach is heuristic. Learning paradigms end-to-end (e.g., via gradient-based method) might discover more diverse or expressive solving patterns.

\textbf{Multi-Solver Integration:} PRISM currently uses Z3 as the sole verification oracle. Integrating multiple solvers (MiniZinc, CPLEX, SMT solvers for different theories) could enable more diverse reasoning strategies and cross-solver knowledge transfer.

\textbf{Interactive Paradigm Refinement:} Offline paradigm libraries are static. Mechanisms for human-in-the-loop refinement or online adaptation could further improve paradigm quality.

\subsection{Broader Implications}

PRISM demonstrates that formal verifiability enables more reliable experience accumulation than natural language heuristics. This insight extends beyond CSP solving:

\begin{itemize}
    \item \textbf{Formal Methods + ML:} Any domain where reasoning steps can be automatically verified (theorem proving, program synthesis, mathematical reasoning) could benefit from formally-grounded memory.

    \item \textbf{Neuro-Symbolic Integration:} Rather than treating neural and symbolic components as separate, viewing the symbolic component (solver, verifier) as a \textit{quality oracle} for neural experiences opens new avenues for experience accumulation in hybrid systems.

    \item \textbf{Repair Strategies:} The four-level escalation protocol is domain-agnostic and could apply to other domains with multiple repair mechanisms (code generation with testing, dialogue with self-evaluation, etc.).
\end{itemize}

\subsection{Final Remarks}

PRISM advances the state of neuro-symbolic reasoning by showing that formally-grounded memory, combined with adaptive repair strategies, can substantially improve LLM performance on structured reasoning tasks. The work is a step toward AI systems that not only learn from experience but verify and integrate that experience reliably.

\section*{Acknowledgments}

We acknowledge support from [institutional funding]. We thank [colleagues] for feedback on early drafts.
